\section{Introduction}

Recent advances in Large Language Models (LLMs) \citep{brown2020language, chung2024scaling} have led to the emergence of search agents \citep{search-o1, search-r1, r1-searcher, deepresearcher, ReSearch} that navigate complex information environments through iterative planning and tool use \citep{react}. Unlike conventional single-step retrieval systems, these agents actively formulate search queries, inspect retrieved observations, and adapt subsequent actions to answer complex, multi-faceted questions.
Reinforcement Learning (RL) has become a standard framework for optimizing such sequential decision-making processes \citep{ppo, grpo}. However, existing RL-based search agents typically rely on gold supervision, such as ground-truth answers, to define reward signals. This reliance creates a fundamental scalability bottleneck: in specialized or rapidly evolving domains, such supervision is often prohibitively expensive or unavailable, making it difficult to construct reliable rewards for search trajectories.

To address this limitation, we propose \textbf{Cycle-Consistent Search (CCS)}, a gold-supervision-free framework for training search agents using the internal structure of the search process itself. Our approach is inspired by cycle-consistency techniques from unsupervised machine translation \citep{he2016dual, lample2017unsupervised, han2021unsupervised} and image-to-image translation \citep{cyclegan, liu2017unsupervised, huang2018multimodal}, where the quality of a transformation is assessed by whether the original input can be recovered through an inverse mapping.
We extend this idea to search by treating a search trajectory $\tau$ as an information-preserving encoding of the source question $q$. As illustrated in \Cref{fig:ccs_ex}, trajectories that omit essential steps for multi-hop questions (middle) or drift toward irrelevant information (right) fail to preserve the information needed to recover the original question, whereas high-quality trajectories (left) do. Formally, our central hypothesis is that, under a high-quality policy $\pi$, the source question should be recoverable from the resulting trajectory: $q \xrightarrow{\pi} \tau \xrightarrow{\phi} \hat{q} \approx q$, where $\phi$ denotes a reconstruction function that maps a search trajectory back to the source question. Under this view, reconstruction quality provides a proxy for trajectory quality.

A central challenge in applying cycle-consistency to search is \emph{information leakage}. Search trajectories may contain superficial lexical cues that enable reconstruction without reflecting genuine search quality. For example, search queries often overlap with the wording of the original question, and final responses may partially restate the question for context, allowing the original question to be reconstructed even when the underlying search process is uninformative or low quality. To mitigate this issue, we introduce two information bottlenecks: excluding the final response from the trajectory and applying Named Entity Recognition (NER) masking \citep{yamada2020luke, sun2019ernie, liu2020k} to search queries, replacing entities such as names and locations with generic tags (e.g., \texttt{[LOC]}). These constraints reduce direct lexical shortcuts and encourage reconstruction to depend on retrieved observations together with the structural scaffold (\Cref{fig:ccs_main}).

\begin{figure*}[!t]
    \centering
    \includegraphics[width=1\linewidth]{img/ccs_ex.pdf}
    \caption{\textbf{Examples of search trajectories.} A high-quality trajectory (left) constitutes a lossless encoding of the question's intent, preserving all necessary information required to reconstruct the original question. In contrast, trajectories that involve insufficient (middle) or irrelevant (right) search steps fail to preserve this information, resulting in incomplete or distorted representations of the question.}
    \label{fig:ccs_ex}
\end{figure*}

Under this formulation, training a search agent amounts to enforcing cycle-consistency over search trajectories, with reconstruction quality on an information-bottlenecked trajectory $\tilde{\tau}$ serving as a proxy reward for trajectory quality. We evaluate CCS on seven representative question-answering benchmarks and a deep research benchmark. The results show that CCS achieves performance comparable to supervised baselines while outperforming existing methods that do not rely on gold supervision. These findings suggest that cycle-consistency can provide a scalable proxy objective for training search agents in settings where gold supervision is unavailable.

In summary, our contributions are as follows:
\begin{itemize}
    \item We introduce Cycle-Consistent Search (CCS), a gold-supervision-free framework for training search agents by enforcing cycle-consistency between questions and search trajectories.
    \item We identify information leakage as a practical challenge in cycle-consistent search training and use information bottlenecks to reduce lexical shortcuts.
    \item We show empirically, across seven question-answering benchmarks and a deep research benchmark, that CCS achieves performance competitive with supervised methods while outperforming prior approaches that do not rely on gold supervision.
\end{itemize}